\chapter{Fundamenta\c{c}\~ao Te\'orica e Trabalhos Relacionados}
\label{cap:referencial}

Este cap\'itulo apresenta os principais conceitos que sustentam o desenvolvimento deste trabalho. Como a proposta envolve o uso de dados escolares para prever desempenho acad\^emico e apoiar decis\~oes pedag\'ogicas, a fundamenta\c{c}\~ao foi organizada em torno da Minera\c{c}\~ao de Dados Educacionais, do Learning Analytics, da previs\~ao de desempenho escolar, da modelagem temporal, da valida\c{c}\~ao dos modelos e dos cuidados \'eticos no uso de dados educacionais.

Tamb\'em s\~ao discutidos trabalhos relacionados que ajudam a situar a proposta desta monografia em um campo de pesquisa j\'a consolidado, mas ainda em expans\~ao, especialmente quando se considera o uso pr\'atico desses modelos no cotidiano das institui\c{c}\~oes de ensino.

\section{Minera\c{c}\~ao de Dados Educacionais e Learning Analytics}
\label{sec:edm-la}

A Minera\c{c}\~ao de Dados Educacionais (EDM) dedica-se ao desenvolvimento de m\'etodos para explorar dados oriundos de contextos educacionais com o objetivo de compreender, monitorar e melhorar processos de aprendizagem \cite{romero2010,handbookedm}. De forma complementar, Learning Analytics (LA) enfatiza a medi\c{c}\~ao, coleta, an\'alise e relato de dados sobre aprendizes e seus contextos para apoiar decis\~oes pedag\'ogicas e institucionais \cite{siemens2012,baker2014}.

Neste trabalho, essas duas perspectivas aparecem de forma integrada. Do lado de EDM, o projeto constr\'oi atributos, testa modelos supervisionados e compara m\'etricas. Do lado de LA, o foco est\'a em transformar resultados anal\'iticos em sa\'idas compreens\'iveis para a rotina escolar. Essa combina\c{c}\~ao \'e importante porque o problema estudado n\~ao termina na previs\~ao em si: ele s\'o faz sentido quando a sa\'ida pode ser lida e usada por quem toma decis\~oes na escola.

\section{Previs\~ao de desempenho escolar}
\label{sec:previsao-desempenho}

Prever desempenho escolar pode assumir diferentes formas: prever evas\~ao, reprova\c{c}\~ao, engajamento, necessidade de refor\c{c}o ou nota futura. Este trabalho escolhe como foco principal a previs\~ao da pr\'oxima nota, porque esse alvo preserva maior granularidade pedag\'ogica e permite interven\c{c}\~oes mais precoces. Em paralelo, o projeto treina um classificador pr\'oprio para estimar o risco de a pr\'oxima nota ficar abaixo de 6,0, aproximando a modelagem do uso institucional sem reduzir o alerta a uma regra simples aplicada sobre a nota prevista.

\section{Modelos de aprendizado de m\'aquina no projeto}
\label{sec:modelos-projeto}

O projeto avaliou modelos supervisionados adequados a dados tabulares, com hist\'orico irregular entre alunos, disciplinas e anos letivos. A escolha dos candidatos levou em conta quatro crit\'erios pr\'aticos: capacidade de capturar rela\c{c}\~oes n\~ao lineares, robustez diante de atributos heterog\^eneos, custo razo\'avel de treinamento e possibilidade de comparar o desempenho com regras simples de refer\^encia. Assim, a discuss\~ao sobre modelos neste cap\'itulo n\~ao tem apenas car\'ater conceitual: ela ajuda a explicar por que certos caminhos foram testados, descartados ou mantidos na vers\~ao final do projeto.

Algoritmos como regress\~ao linear e \'arvore de decis\~ao simples seriam mais f\'aceis de explicar, mas tenderiam a representar de forma limitada intera\c{c}\~oes entre nota anterior, tend\^encia, faltas, pagamentos e contexto escolar. M\'etodos como Support Vector Machine (SVM) poderiam ser usados, mas costumam exigir maior cuidado com escala, ajuste de hiperpar\^ametros e interpreta\c{c}\~ao em bases tabulares maiores. Redes neurais profundas, por sua vez, foram consideradas menos aderentes ao volume e \`a irregularidade das sequ\^encias dispon\'iveis. Modelos como CatBoost e LightGBM permanecem como alternativas futuras, mas n\~ao foram priorizados na vers\~ao final para manter a solu\c{c}\~ao mais simples, reproduz\'ivel e apoiada em bibliotecas j\'a consolidadas no ambiente do projeto. Em outras palavras, a sele\c{c}\~ao dos candidatos buscou equilibrar desempenho, legibilidade e viabilidade real de manuten\c{c}\~ao.

\subsection{Random Forest}
\label{subsec:random-forest}

O Random Forest, proposto por Breiman \cite{breiman2001}, foi importante no projeto por servir como primeiro modelo tabular forte ap\'os a fase explorat\'oria. Sua principal vantagem est\'a na robustez em dados tabulares, na capacidade de capturar rela\c{c}\~oes n\~ao lineares e na boa toler\^ancia a atributos heterog\^eneos. No desenvolvimento deste trabalho, ele teve um papel de refer\^encia concreta: ajudou a estabelecer um patamar forte de compara\c{c}\~ao antes da consolida\c{c}\~ao de outros candidatos.

\subsection{LSTM}
\label{subsec:lstm}

As redes Long Short-Term Memory (LSTM), introduzidas por Hochreiter e Schmidhuber \cite{hochreiter1997}, foram consideradas por causa da natureza temporal do problema. No entanto, o conjunto de dados sintéticos adotado na vers\~ao p\'ublica preservou sequ\^encias irregulares, entradas e sa\'idas de alunos e hist\'oricos de extens\~ao muito desigual para manter a complexidade estrutural do problema. Esses fatores reduziram a ader\^encia dessa abordagem ao escopo institucional do trabalho. O ponto aqui n\~ao \'e afirmar que LSTM seja inadequada em termos absolutos, mas reconhecer que, neste conjunto de dados e neste problema, ela n\~ao parecia ser a escolha mais consistente.

\subsection{Gradient Boosting}
\label{subsec:gradient-boosting}

O Gradient Boosting, formalizado por Friedman \cite{friedman2001}, mostrou-se especialmente adequado \`a solu\c{c}\~ao final por trabalhar bem com dados tabulares e atributos derivados. A varia\c{c}\~ao HistGradientBoosting foi inclu\'ida por ser eficiente em bases tabulares e por lidar bem com padr\~oes n\~ao lineares. Na rodada final, ele foi o melhor candidato para o modo de alerta de risco, enquanto o Random Forest permaneceu competitivo e foi selecionado para a regress\~ao no modo de previs\~ao de nota. Esse resultado refor\c{c}a que o projeto n\~ao convergiu para um \'unico modelo universal, mas para escolhas diferentes conforme a natureza de cada tarefa.

\subsection{Classifica\c{c}\~ao supervisionada e regras simples}
\label{subsec:classificacao-regras}

Uma decis\~ao importante foi manter a classifica\c{c}\~ao de risco como tarefa supervisionada, em vez de utilizar apenas regras fixas, como ``se a nota prevista for menor que 6, ent\~ao o aluno est\'a em risco''. Essa regra seria simples, mas perderia parte do comportamento temporal do problema. O aluno pode ter uma nota prevista numericamente pr\'oxima da m\'edia e, ainda assim, apresentar padr\~ao de queda, hist\'orico anterior fraco na disciplina, faltas acumuladas ou dist\^ancia negativa em rela\c{c}\~ao \`a turma. A classifica\c{c}\~ao aprende diretamente com exemplos hist\'oricos quais combina\c{c}\~oes desses sinais antecederam notas futuras abaixo de 6,0. Essa escolha ajuda a aproximar o alerta da din\^amica real do acompanhamento escolar.

Isso n\~ao elimina as regras. No projeto, regras e pontua\c{c}\~oes continuam sendo usadas na camada de relat\'orios e prioriza\c{c}\~ao. A diferen\c{c}a \'e que o alerta principal de risco \'e estimado por um classificador treinado sobre dados sintéticos compat\'iveis com o contrato p\'ublico, enquanto as regras ajudam a organizar a comunica\c{c}\~ao pedag\'ogica e a calibrar prioridades.

\section{Dados temporais, validade e m\'etricas}
\label{sec:dados-temporais-metricas}

Um ponto cr\'itico em aplica\c{c}\~oes educacionais \'e a valida\c{c}\~ao temporal. Em vez de dividir aleatoriamente treino e teste, este trabalho adota separa\c{c}\~ao por anos letivos, evitando vazamento entre passado e futuro. Essa decis\~ao aproxima a avalia\c{c}\~ao do uso real da solu\c{c}\~ao, em que a escola precisa prever o pr\'oximo per\'iodo a partir do passado j\'a conhecido. A preocupa\c{c}\~ao com valida\c{c}\~ao temporal tamb\'em \'e discutida em estudos sobre avalia\c{c}\~ao de modelos em dados dependentes do tempo, nos quais uma divis\~ao aleat\'oria pode produzir estimativas otimistas de desempenho \cite{bergmeir2018}. Para este TCC, esse cuidado \'e central porque o objetivo n\~ao era apenas obter uma m\'etrica melhor, mas testar uma situa\c{c}\~ao mais parecida com o uso futuro da pipeline.

As m\'etricas escolhidas acompanham essa dupla natureza do problema. Para regress\~ao, o trabalho utiliza Mean Absolute Error (MAE), que mede o erro m\'edio absoluto em pontos de nota, e Root Mean Squared Error (RMSE), que penaliza erros maiores com mais intensidade. Para classifica\c{c}\~ao de risco, utiliza precision, recall e F1-score (F1). A precision indica, entre os alunos sinalizados, quantos realmente estavam em risco; o recall indica, entre os alunos que realmente estavam em risco, quantos foram recuperados pelo modelo; e o F1 combina essas duas medidas. Essa combina\c{c}\~ao permite avaliar tanto proximidade num\'erica quanto capacidade de recuperar casos cr\'iticos.

\section{\'Etica, privacidade e interpretabilidade}
\label{sec:etica-privacidade-interpretabilidade}

O uso de dados educacionais exige aten\c{c}\~ao especial \`a privacidade, \`a finalidade do tratamento e \`a interpreta\c{c}\~ao cuidadosa das sa\'idas. Discuss\~oes em LA e EDM mostram que prote\c{c}\~ao de dados e governan\c{c}a n\~ao s\~ao barreiras para an\'alise, mas condi\c{c}\~oes para seu uso institucionalmente leg\'itimo \cite{slade2013,gasevic2016,khalil2016}. No contexto brasileiro, a Lei Geral de Prote\c{c}\~ao de Dados Pessoais (LGPD) refor\c{c}a esse cuidado \cite{lgpd}. No caso deste projeto, essa discuss\~ao n\~ao aparece apenas como requisito formal: ela interfere diretamente na forma de extrair, preparar, armazenar e comunicar os resultados.

Por isso, a vers\~ao publicada do projeto passou a trabalhar com dados sintéticos e assumiu, desde o in\'icio, que a sa\'ida do sistema n\~ao poderia ser um veredito definitivo sobre o aluno. Os relat\'orios foram desenhados para indicar fatores de alerta e recomenda\c{c}\~oes iniciais, preservando o papel humano na decis\~ao pedag\'ogica \cite{fiok2022}.

\section{Trabalhos relacionados}
\label{sec:trabalhos-relacionados}

Os trabalhos relacionados ao tema podem ser agrupados em tr\^es eixos. O primeiro re\'une revis\~oes gerais da \'area, como as de Romero e Ventura \cite{romero2010} e de Baker e Inventado \cite{baker2014}. O segundo eixo envolve estudos preditivos aplicados \`a educa\c{c}\~ao, sobretudo aqueles que combinam m\'ultiplas fontes de informa\c{c}\~ao. O terceiro eixo aborda interpretabilidade, governan\c{c}a e quest\~oes \'eticas. Entre esses conjuntos, os trabalhos mais diretamente ligados a esta monografia s\~ao os que tratam de previs\~ao de nota e detec\c{c}\~ao de risco acad\^emico.

\subsection{Compara\c{c}\~ao com estudos de previs\~ao de nota e detec\c{c}\~ao de risco}
\label{subsec:comparacao-trabalhos}

Considerando o foco deste trabalho em prever a pr\'oxima nota e sinalizar risco pedag\'ogico, a literatura mais pr\'oxima pode ser organizada em dois subconjuntos: estudos voltados \`a previs\~ao de nota futura e estudos voltados \`a detec\c{c}\~ao de alunos em risco acad\^emico. No primeiro grupo, Morsy e Karypis \cite{morsy2017} e Widjaja et al. \cite{widjaja2020} ajudam a situar o problema da regress\~ao da pr\'oxima nota dentro de uma linha de pesquisa j\'a consolidada, embora majoritariamente concentrada em ensino superior.

No segundo grupo, Lakkaraju et al. \cite{lakkaraju2015} e Koprinska, Stretton e Yacef \cite{koprinska2015} representam estudos de detec\c{c}\~ao de risco. Nesse eixo, a proximidade com esta pesquisa aparece menos no tipo exato de alvo e mais na preocupa\c{c}\~ao com identifica\c{c}\~ao precoce, interpretabilidade e utilidade institucional. As Tabelas~\ref{tab:trab-rel-prev-nota} e~\ref{tab:trab-rel-risco} sintetizam esses dois recortes comparativos.

\begin{table}[!htbp]
\caption{Trabalhos correlatos focados em previs\~ao de nota futura.}
\label{tab:trab-rel-prev-nota}
\centering
\footnotesize
\begin{tabularx}{\textwidth}{L{3.0cm} X X X}
\hline
\textbf{Trabalho} & \textbf{Entradas} & \textbf{Sa\'ida e m\'etodos} & \textbf{Limite frente a este TCC} \\
\hline
Morsy e Karypis (2017) & Hist\'orico aluno-disciplina e rela\c{c}\~oes de conhecimento entre cursos em base universit\'aria & Regress\~ao da nota em disciplinas do pr\'oximo per\'iodo por meio de \textit{Cumulative Knowledge-based Regression Models} & Foco em matr\'icula universit\'aria futura; n\~ao trata alerta pedag\'ogico por etapa nem relat\'orios operacionais escolares \\
Widjaja et al. (2020) & Hist\'orico aluno-curso, atributos de aluno e atributos de curso em contexto universit\'ario & Predi\c{c}\~ao de nota do pr\'oximo termo com \textit{Factorization Machine} e LSTM-FM & Foco em desempenho por termo e escolha de cursos; n\~ao integra risco pedag\'ogico separado nem sinais administrativos escolares \\
\hline
\end{tabularx}
\fontetab{Elaborada pelo autor com base nos trabalhos citados.}
\end{table}

\begin{table}[!htbp]
\caption{Trabalhos correlatos focados em detec\c{c}\~ao de risco acad\^emico.}
\label{tab:trab-rel-risco}
\centering
\footnotesize
\begin{tabularx}{\textwidth}{L{3.1cm} X X X}
\hline
\textbf{Trabalho} & \textbf{Entradas} & \textbf{Sa\'ida e m\'etodos} & \textbf{Limite frente a este TCC} \\
\hline
Lakkaraju et al. (2015) & Hist\'orico escolar, notas, aus\^encias, indicadores demogr\'aficos e administrativos em dois distritos dos EUA & Classifica\c{c}\~ao bin\'aria de risco de n\~ao concluir o ensino m\'edio no tempo esperado, com compara\c{c}\~ao entre classificadores e visualiza\c{c}\~ao para gestores & O alvo \'e conclus\~ao escolar de longo prazo; n\~ao prev\^e a pr\'oxima nota da disciplina nem trabalha com sa\'ida combinada de regress\~ao e risco \\
Koprinska, Stretton e Yacef (2015) & Notas parciais, submiss\~oes em sistema de corre\c{c}\~ao autom\'atica e atividade em f\'orum de discuss\~ao em disciplina de programa\c{c}\~ao & Classifica\c{c}\~ao de aprova\c{c}\~ao ou falha no exame final com \'arvore de decis\~ao e regras acion\'aveis para remedia\c{c}\~ao & Estudo concentrado em uma disciplina universit\'aria e em dados de ambiente digital; n\~ao integra contexto financeiro, professor, s\'erie ou relat\'orios multiatores \\
\hline
\end{tabularx}
\fontetab{Elaborada pelo autor com base nos trabalhos citados.}
\end{table}

\FloatBarrier

Esses estudos mostram que a previs\~ao de nota futura j\'a est\'a consolidada como problema relevante, sobretudo em ensino superior. Entretanto, o alvo adotado costuma ser a nota em disciplinas ainda n\~ao cursadas ou em um pr\'oximo termo acad\^emico. Aqui, a sa\'ida de regress\~ao foi reposicionada para a pr\'oxima nota observada dentro da sequ\^encia temporal da mesma disciplina, o que favorece interven\c{c}\~oes pedag\'ogicas de curto prazo no cotidiano escolar. Essa mudan\c{c}a de foco \'e pequena em apar\^encia, mas importante em termos de uso pr\'atico.

Nos estudos de risco, nota-se forte preocupa\c{c}\~ao com identifica\c{c}\~ao precoce, interpretabilidade e apoio \`a interven\c{c}\~ao. Esse ponto aproxima-se diretamente deste TCC. A diferen\c{c}a principal est\'a no fato de que o problema aqui n\~ao trata risco apenas como evas\~ao, conclus\~ao tardia ou aprova\c{c}\~ao final em uma disciplina. O alvo bin\'ario foi definido como o risco de a pr\'oxima nota observada ficar abaixo de 6,0, o que torna o alerta mais pr\'oximo da linguagem pedag\'ogica da escola e mais aderente a decis\~oes de curto prazo.

Do ponto de vista metodol\'ogico, o trabalho combina ideias que na literatura frequentemente aparecem separadas. Dos estudos de previs\~ao de nota, aproveita a l\'ogica de usar o hist\'orico acad\^emico como base para regress\~ao temporal. Dos estudos de risco, incorpora a necessidade de triagem antecipada, prioriza\c{c}\~ao de casos e utilidade institucional. A diferen\c{c}a \'e que essas duas perspectivas foram unificadas em uma mesma arquitetura temporal com duas sa\'idas supervisionadas: \texttt{target\_nota\_proxima} e \texttt{target\_risco\_proxima}.

Tamb\'em h\'a diferen\c{c}as importantes nas entradas e nas ferramentas anal\'iticas. Enquanto os trabalhos universit\'arios de previs\~ao de nota concentram-se sobretudo em rela\c{c}\~oes aluno-curso e atributos acad\^emicos, e os trabalhos de risco costumam enfatizar desempenho, engajamento ou indicadores administrativos amplos, este TCC articula notas, faltas, pagamentos, professor e contexto escolar em uma mesma base temporal. Em termos de modelagem, utiliza Random Forest, HistGradientBoosting e baselines heur\'isticos como refer\^encia honesta de compara\c{c}\~ao, todos avaliados por valida\c{c}\~ao temporal por ano letivo. Em termos de sa\'ida, a solu\c{c}\~ao n\~ao termina na m\'etrica: as predi\c{c}\~oes alimentam relat\'orios para professor, coordena\c{c}\~ao e secretaria, al\'em de ranking executivo de prioridade.

Em rela\c{c}\~ao a esses trabalhos, este TCC se diferencia por articular toda a trajet\'oria de um projeto aplicado: tratamento do banco, refatora\c{c}\~ao arquitetural, valida\c{c}\~ao temporal, gera\c{c}\~ao de relat\'orios e ranking executivo. Em vez de comparar apenas algoritmos em um dataset pronto, o trabalho documenta a constru\c{c}\~ao completa de uma solu\c{c}\~ao voltada ao contexto escolar, reunindo previs\~ao da pr\'oxima nota, alerta de risco pedag\'ogico e operacionaliza\c{c}\~ao institucional em uma mesma pipeline reproduz\'ivel. Esse posicionamento ajuda a explicar a contribui\c{c}\~ao do trabalho n\~ao apenas como experimento de modelagem, mas como solu\c{c}\~ao anal\'itica aplicada.
